\chapter{Light, bands and the hyperspectral cube}\label{ch:light}

\begin{onething}
A hyperspectral pixel is not a colour, it is a curve: one reflectance value per narrow band of
wavelengths, which you can treat as a vector of $C$ numbers. That curve identifies what the
pixel is made of, and a whole image is a three-dimensional array of such curves.
\end{onething}

\section{Why three numbers are not enough}\label{sec:light:motivation}

A digital photograph stores three numbers per pixel: how much red, how much green and how much
blue. Your eye works the same way, with three kinds of colour receptor, so three numbers are
enough to fool it. Three numbers are also very little information. Fresh grass and a sheet of
green-painted metal can produce the same red, green and blue triple, and once the three numbers
match, nothing distinguishes them.

The sensors in this book record \emph{hundreds} of numbers per pixel instead of three, most of
them for light your eye cannot see. With those extra numbers grass and green paint stop looking
alike, because grass has a sharp rise in reflected light immediately beyond red and paint does not. This
chapter builds the vocabulary for that statement, carefully and with arithmetic: what a
wavelength is, what ``reflectance'' measures, what a band actually records, and how all of it is
stored as an array.

\section{Light and wavelength}\label{sec:light:wavelength}

Light is an electromagnetic wave. Its \textbf{wavelength}\index{wavelength}, written $\lambda$
(the Greek letter ``lambda''), is the distance from one crest of the wave to the next. Different
wavelengths are what we perceive as different colours, and beyond the ends of the visible range
the same physics continues invisibly.

Wavelengths of light are tiny, so they are measured in \textbf{nanometres}\index{nanometre}
(nm):
\[ 1\ \text{nm} = 10^{-9}\ \text{m} = 0.000\,000\,001\ \text{m}. \]
So $1{,}650$ nm is $1{,}650 \times 10^{-9}$ m $= 1.65 \times 10^{-6}$ m, which is $1.65$
micrometres. Remote-sensing papers often quote micrometres ($\mu$m); this codebase uses
nanometres everywhere.

\begin{table}[htbp]\centering
\caption{The wavelength ranges used in this book. Boundaries between the named regions are
conventions, not physics.}\label{tab:light:ranges}
\begin{tabular}{@{}lll@{}}\toprule
Region & Range (nm) & What it shows \\\midrule
Visible (VIS) & about 400--700 & colour as you see it\\
Near infrared (NIR) & about 700--1{,}000 & leaf structure, the ``red edge''\\
Shortwave infrared (SWIR) & about 1{,}000--2{,}500 & water, minerals, clays\\
Thermal infrared (TIR) & about 10{,}000--12{,}000 & emitted heat (temperature)\\\bottomrule
\end{tabular}
\end{table}

The hyperspectral sensors in this book cover roughly $400$ to $2{,}500$ nm and the shared grid of
\cref{ch:sensors} keeps $460$ to $2{,}450$ nm. Most of that is invisible. Two of the five sensors
are thermal: they record emitted heat near $10{,}000$ nm and are treated separately in
\cref{sec:light:thermal}.

\section{Radiance, irradiance and reflectance}\label{sec:light:reflectance}

A satellite does not measure ``what the ground is''. It measures light arriving at its detector.
To turn that into a property of the ground you need three quantities, which are easy to confuse
because all three are ``amounts of light''.

\begin{definition}[Irradiance]\index{irradiance}
\textbf{Irradiance} $E(\lambda)$ is the power of sunlight \emph{arriving} on a surface, per
square metre of surface and per nanometre of wavelength. Units: W\,m$^{-2}$\,nm$^{-1}$.
\end{definition}

\begin{definition}[Radiance]\index{radiance}
\textbf{Radiance} $L(\lambda)$ is the power \emph{leaving} a surface in one particular direction,
per square metre, per unit of solid angle (a ``steradian'', sr, the three-dimensional version of
an angle), and per nanometre. Units: W\,m$^{-2}$\,sr$^{-1}$\,nm$^{-1}$. A sensor looking down at
the ground measures radiance.
\end{definition}

\begin{definition}[Reflectance]\index{reflectance}
\textbf{Reflectance} $\rho(\lambda)$ (Greek ``rho'') is the fraction of the arriving light that
the surface sends back, at wavelength $\lambda$. Because it is a ratio of two amounts of light,
it has no units.
\end{definition}

\subsection{Reflectance as a ratio}

In its simplest form reflectance is ``what came back divided by what arrived''. If $1.20$ units of
light arrive at some wavelength and $0.30$ units come back,
\[ \rho = \frac{0.30}{1.20} = 0.25, \]
so the surface reflects a quarter of the light at that wavelength. A value of $0$ would mean the
surface absorbs everything (perfectly black); $1$ would mean it returns everything.

The ratio is what makes reflectance useful. If the sun is lower, both the arriving and the
returned light fall, and their ratio stays the same. If a different satellite measures the same
field, its electronics differ but the surface's ratio does not. Reflectance is a property of the
\emph{surface}, which is what lets one model work across sensors and lets a pixel be compared with
a laboratory measurement of a known material.

\subsection{From radiance to reflectance: where the \texorpdfstring{$\pi$}{pi} comes from}

The sensor measures radiance in one direction, but the surface sends light in all directions.
To compare ``what came back'' with ``what arrived'' you need the total returned over the whole sky
above the surface. The standard simplifying assumption is a
\textbf{Lambertian}\index{Lambertian surface} surface: one that looks equally bright from every
viewing direction, so its radiance $L$ is the same in all directions. Matt surfaces such as soil
are close to this; water and metal are not.

The total power per square metre leaving such a surface, called its exitance $M$, adds up the
radiance over every direction of the upper hemisphere, weighting each direction by
$\cos\theta$ (a beam leaving at angle $\theta$ from vertical crosses a horizontal square metre
at a slant, so only $\cos\theta$ of it counts). Measure directions by $\theta$ from vertical
($0$ to $\pi/2$) and $\varphi$ around the vertical ($0$ to $2\pi$); a small patch of directions
has solid angle $\sin\theta\,d\theta\,d\varphi$. Then:
\begin{align}
M &= \int_{0}^{2\pi}\!\!\int_{0}^{\pi/2} L \cos\theta \sin\theta \, d\theta\, d\varphi.
\intertext{$L$ is the same in every direction (Lambertian), so take it outside both integrals:}
M &= L \int_{0}^{2\pi} d\varphi \int_{0}^{\pi/2} \cos\theta\sin\theta\, d\theta.
\intertext{The $\varphi$ integral is the length of the range, $2\pi$. For the $\theta$ integral
note that the derivative of $\tfrac{1}{2}\sin^2\theta$ is $\sin\theta\cos\theta$, so}
\int_{0}^{\pi/2}\cos\theta\sin\theta\,d\theta &= \Big[\tfrac{1}{2}\sin^2\theta\Big]_{0}^{\pi/2}
 = \tfrac{1}{2}(1)^2 - \tfrac{1}{2}(0)^2 = \tfrac{1}{2}.
\intertext{Multiply the two pieces:}
M &= L \times 2\pi \times \tfrac{1}{2} = \pi L. \label{eq:light:exitance}
\end{align}
Sunlight arriving from a sun at angle $\theta_s$ from vertical is spread over a larger patch of
ground than when the sun is overhead, so the irradiance on a horizontal surface is
$E\cos\theta_s$, where $E$ is the irradiance on a surface facing the sun. Reflectance is returned
over arrived:
\begin{equation}\label{eq:light:rho}
\rho = \frac{M}{E\cos\theta_s} = \frac{\pi L}{E \cos\theta_s}.
\end{equation}

\begin{worked}[label={we:light:rho}]{Radiance to reflectance}
The numbers here are illustrative\tagA, chosen to be of a realistic size. At $860$ nm a sensor
measures radiance $L = 0.10$ W\,m$^{-2}$\,sr$^{-1}$\,nm$^{-1}$; the solar irradiance is
$E = 1.50$ W\,m$^{-2}$\,nm$^{-1}$ and the sun is $\theta_s = 30^\circ$ from vertical. Do each step
on paper.

The cosine of the sun angle:
\[ \cos 30^\circ \approx 0.8660. \]
Numerator of \cref{eq:light:rho}:
\[ \pi L = 3.1416 \times 0.10 = 0.31416. \]
Denominator:
\[ E\cos\theta_s = 1.50 \times 0.8660 = 1.2990. \]
Divide:
\[ \rho = 0.31416 / 1.2990 \approx 0.242. \]
The surface returns about $24\%$ of the light at $860$ nm, typical of vegetation in the near
infrared. If the sun dropped to $60^\circ$ ($\cos 60^\circ = 0.5$) and the surface were unchanged,
the arriving light would fall to $1.50 \times 0.5 = 0.75$ and the measured radiance would fall in
proportion to $L = 0.242 \times 0.75 / \pi \approx 0.0578$; the ratio $\rho$ would still come out
as $0.242$.
\end{worked}

This is background physics. In this codebase the satellite providers have already done the
correction for the atmosphere and the sun: PRISMA's \emph{L2D} and EnMAP's \emph{L2A} products
are surface reflectance stored as integers, and the code's job is only to undo the integer
scaling (\cref{ch:sensors}). The ``L2'' prefix means ``level 2'', a processing level at which the
provider delivers a physical surface quantity rather than raw counts.

\subsection{What real reflectance looks like}

Real ground sits mostly between about $0.02$ and $0.30$: most surfaces return between $2\%$ and
$30\%$ of the light. You can check this against the project's own data. The file
\file{app/constants/hyperspectral.json} holds the average of every band over the training corpus
(\cref{ch:bandpipeline} uses it). Its first value, for $460$ nm, is
$0.04301$\tagM{app/constants/hyperspectral.json}, and the largest of its $165$ averages is $0.2775$
at $1{,}270$ nm\tagC{} (computed by \file{research/textbook/tools/compute_light.py}). Both are in
the range quoted above.

Values above $1$ are not automatically bugs. Sunlight glinting off water or metal can be brighter
than a perfect matt reflector would give, and after atmospheric correction it can exceed $1$.
Do not add a hard clamp at $1$ assuming otherwise.

\section{What a band records}\label{sec:light:band}

\begin{definition}[Band]\index{band}
A \textbf{band} is one measurement channel of a sensor: the amount of light it records in one
narrow range of wavelengths. ``Red'' is a band; so is ``light between $1{,}645$ and $1{,}655$ nm''.
In deep-learning code the same thing is called a \textbf{channel}\index{channel}; this codebase
uses both words, often in the same file.
\end{definition}

A band is not a knife-edge at one wavelength. A detector element responds to a small range of
wavelengths, most strongly at its centre and less towards the sides. The shape of that response
is the band's \textbf{spectral response function}\index{spectral response function} (SRF), a
curve $w(\lambda) \ge 0$. What the band reports is a \emph{weighted average} of the scene's
spectrum, with the SRF as the weights:
\begin{equation}\label{eq:light:srf}
b = \frac{\int w(\lambda)\,\rho(\lambda)\,d\lambda}{\int w(\lambda)\,d\lambda}
\;\approx\; \frac{\sum_k w_k\,\rho(\lambda_k)}{\sum_k w_k}.
\end{equation}
The right-hand form samples the curve at a few wavelengths $\lambda_k$ with weights
$w_k = w(\lambda_k)$. Read it in words: multiply each reflectance by how strongly the band
responds at that wavelength, add the products, and divide by the total weight so that a flat
spectrum comes back unchanged.

\begin{worked}[label={we:light:srf}]{A band as a weighted average of three samples}
A band is centred at $1{,}650$ nm. Sample its response at three wavelengths: $1{,}645$, $1{,}650$
and $1{,}655$ nm, with weights $0.5$, $1$, $0.5$ (half as sensitive $5$ nm either side of the
centre; \cref{we:light:fwhm} shows where those halves come from). The surface's reflectance at the
three wavelengths is $0.20$, $0.22$, $0.30$ (illustrative\tagA). Compute the band value.

Multiply each reflectance by its weight:
\begin{align*}
0.5 \times 0.20 &= 0.10,\\
1 \times 0.22 &= 0.22,\\
0.5 \times 0.30 &= 0.15.
\end{align*}
Add the products, and add the weights:
\begin{align*}
0.10 + 0.22 + 0.15 &= 0.47,\\
0.5 + 1 + 0.5 &= 2.
\end{align*}
Divide:
\[ b = 0.47 / 2 = 0.235. \]
Compare three possible answers: the reflectance at the centre alone is $0.22$; the plain average of
the three samples is $(0.20 + 0.22 + 0.30)/3 = 0.72/3 = 0.24$; the band reports $0.235$. The band
is pulled towards the steeply rising side of the spectrum, but less than an unweighted average
would be. Check the normalisation: if all three reflectances were $0.22$, the numerator would be
$0.5\times0.22 + 0.22 + 0.5\times0.22 = 0.44$ and the band would report $0.44/2 = 0.22$, unchanged.
\end{worked}

\subsection{Band width: the full width at half maximum}

A common model of an SRF is the bell-shaped \textbf{Gaussian} curve
\begin{equation}\label{eq:light:gauss}
w(\lambda) = \exp\!\left(-\frac{(\lambda - \lambda_c)^2}{2\sigma^2}\right),
\end{equation}
centred at the band's \textbf{centre wavelength}\index{centre wavelength} $\lambda_c$, with
width set by $\sigma$. At the centre $w = \exp(0) = 1$. Sensor documentation does not quote
$\sigma$; it quotes the \textbf{full width at half maximum}\index{full width at half maximum}
(FWHM): the distance between the two wavelengths where the response has fallen to half its peak.
The two are related by a fixed factor, which you can derive.

Find the offset $\Delta = \lambda - \lambda_c$ where the response is one half:
\begin{align*}
\exp\!\left(-\frac{\Delta^2}{2\sigma^2}\right) &= \frac{1}{2}.
\intertext{Take the natural logarithm of both sides (the logarithm undoes the exponential):}
-\frac{\Delta^2}{2\sigma^2} &= \ln\tfrac{1}{2} = -\ln 2.
\intertext{Multiply both sides by $-2\sigma^2$:}
\Delta^2 &= 2\sigma^2\ln 2.
\intertext{Take the positive square root:}
\Delta &= \sigma\sqrt{2\ln 2}.
\intertext{The half-maximum points sit at $\lambda_c - \Delta$ and $\lambda_c + \Delta$, so the
full width is twice $\Delta$:}
\text{FWHM} &= 2\sqrt{2\ln 2}\;\sigma.
\end{align*}
Now the number. $\ln 2 \approx 0.6931$, so $2\ln 2 \approx 1.3863$, its square root is about
$1.1774$, and twice that is about $2.3548$:
\begin{equation}\label{eq:light:fwhm}
\text{FWHM} \approx 2.3548\,\sigma, \qquad \sigma \approx \frac{\text{FWHM}}{2.3548}.
\end{equation}

\begin{worked}[label={we:light:fwhm}]{The weights of a 10 nm band}
A band has FWHM $10$ nm. Find $\sigma$, then the Gaussian weight $5$ nm from the centre.

By \cref{eq:light:fwhm}:
\[ \sigma = 10 / 2.3548 \approx 4.2466 \text{ nm}. \]
At $\Delta = 5$ nm, square the offset and the width:
\begin{align*}
\Delta^2 &= 25,\\
2\sigma^2 &= 2 \times 4.2466^2 = 2 \times 18.034 = 36.068.
\end{align*}
Divide and exponentiate:
\begin{align*}
\Delta^2 / (2\sigma^2) &= 25/36.068 \approx 0.6931,\\
w &= \exp(-0.6931) \approx 0.5000.
\end{align*}
Exactly one half, as it must be: $5$ nm is half of the FWHM, and the FWHM is by definition where the
response is half. Those are the weights $0.5$, $1$, $0.5$ used in \cref{we:light:srf}. For PRISMA,
whose bands are about $12$ nm wide, $\sigma \approx 12/2.3548 \approx 5.096$ nm\tagC.
\end{worked}

\begin{figure}[htbp]\centering
\begin{tikzpicture}
\begin{axis}[width=0.9\linewidth,height=5.2cm,xmin=1625,xmax=1675,ymin=0,ymax=1.15,
  xlabel={wavelength $\lambda$ (nm)},ylabel={response $w(\lambda)$},
  axis lines=left,tick label style={font=\small},label style={font=\small},
  xtick={1630,1640,1645,1650,1655,1660,1670}]
\addplot[spec,very thick,domain=1625:1675,samples=120]{exp(-((x-1650)^2)/(2*4.2466^2))};
\addplot[only marks,mark=*,mark size=2pt,loss] coordinates {(1645,0.5) (1650,1) (1655,0.5)};
\draw[dashed,maskgrey] (axis cs:1625,0.5) -- (axis cs:1675,0.5);
\draw[<->,thick,attn] (axis cs:1645,0.42) -- node[below,font=\small]{FWHM $=10$ nm} (axis cs:1655,0.42);
\node[font=\small,anchor=west] at (axis cs:1656,0.95) {$\lambda_c = 1650$ nm, $\sigma \approx 4.25$ nm};
\end{axis}
\end{tikzpicture}
\caption{A Gaussian spectral response function with FWHM $10$ nm, drawn from
\cref{eq:light:gauss}. The response falls to one half exactly $5$ nm either side of the centre;
the three amber points are the samples used in \cref{we:light:srf}.}\label{fig:light:srf}
\end{figure}

Every band therefore carries two numbers of metadata: its centre wavelength and its FWHM. The
codebase stores them as lists, one entry per band (\cref{ch:product}). FWHM matters most for
naming materials: comparing a satellite spectrum with a laboratory spectrum means simulating how
the laboratory curve would look through \emph{this} sensor's bands, which is
\cref{eq:light:srf} applied to the laboratory curve. The companion volume on spectral matching
works that convolution in full. A scene without FWHM values cannot be material-matched at all.

\section{A spectrum is a vector}\label{sec:light:spectrum}

Pick one pixel and read its value in every band in order. You get $C$ numbers. Plotted against
the bands' centre wavelengths they form a curve, the pixel's
\textbf{spectrum}\index{spectrum}. For mathematics it is more useful to treat those $C$ numbers as
a single object.

\begin{definition}[Vector]\index{vector}
A \textbf{vector} in $\R^C$ is an ordered list of $C$ real numbers, written as a column,
\[ \vect{x} = \begin{pmatrix} x_0 \\ x_1 \\ \vdots \\ x_{C-1} \end{pmatrix}, \]
where $x_c$ is its $c$-th \textbf{component}. $\R^C$ (``R to the C'') is the set of all such
lists.
\end{definition}

A pixel's spectrum on the shared grid is a vector $\vect{x} \in \R^{165}$, and the component
$x_c$ is its reflectance in band $c$. Two operations on vectors are used constantly, and both act
component by component.

\begin{itemize}
\item \textbf{Subtraction}: $(\vect{x} - \vect{y})_c = x_c - y_c$. The residual of
\cref{we:problem:residual} was $\vect{x} - \xhat$.
\item \textbf{Scaling} by a number $a$: $(a\vect{x})_c = a\,x_c$.
\end{itemize}

\begin{worked}[label={we:light:vector}]{Two pixels of the same material}
Two pixels on the same field have spectra over three bands
\[ \vect{x} = (0.04,\ 0.24,\ 0.12), \qquad \vect{y} = (0.02,\ 0.12,\ 0.06) \]
(illustrative\tagA; written as rows to save space). The second pixel lies in the shadow of a
hill. Compare them.

Divide each component of $\vect{y}$ by the matching component of $\vect{x}$:
\begin{align*}
0.02/0.04 &= 0.5,\\
0.12/0.24 &= 0.5,\\
0.06/0.12 &= 0.5.
\end{align*}
Every ratio is $0.5$, so $\vect{y} = 0.5\,\vect{x}$: the same shape, half the brightness. Their
difference is
\[ \vect{x} - \vect{y} = (0.04 - 0.02,\ 0.24 - 0.12,\ 0.12 - 0.06) = (0.02,\ 0.12,\ 0.06), \]
which is not small, even though the material is identical. A score built on plain differences
calls the shadowed pixel different; a score built on the \emph{shape} (the direction of the
vector) does not. \Cref{ch:losses} builds that second kind of score, the spectral angle.
\end{worked}

\subsection{What a real spectrum looks like}

\Cref{fig:light:meanspec} plots a real spectrum: the average of each of the $165$ grid bands
over the corpus used to fit Indradhanu's normalisation (\cref{ch:bandpipeline}), with a shaded
band one standard deviation either side. Three features are worth finding on it.

\begin{itemize}
\item \textbf{The red edge}\index{red edge}. Between $680$ and $760$ nm the mean rises from
$0.0777$ to $0.1831$\tagC, more than doubling in $80$ nm. Chlorophyll absorbs red light, and leaf
structure scatters near-infrared light strongly, so vegetated pixels jump sharply here. Nothing
else does that, which is why it separates grass from green paint.
\item \textbf{Water absorption}. The curve is cut at four places, around $940$, $1{,}140$,
$1{,}400$ and $1{,}875$ nm. Water vapour in the air absorbs so strongly there that the bands are
removed from the grid (\cref{ch:sensors}). Next to the gaps the curve dips, for example to
$0.1341$ at $1{,}480$ nm\tagC.
\item \textbf{SWIR fall-off}. Beyond about $1{,}700$ nm reflectance falls, with small dips such as
the one near $2{,}200$ nm where clay minerals absorb.
\end{itemize}

\begin{figure}[htbp]\centering
\begin{tikzpicture}
\begin{axis}[width=\linewidth,height=6.2cm,xmin=440,xmax=2470,ymin=0,ymax=0.40,
  xlabel={wavelength (nm)},ylabel={reflectance},axis lines=left,
  tick label style={font=\small},label style={font=\small},
  xtick={460,700,1000,1300,1600,1900,2200,2450},ytick={0,0.1,0.2,0.3,0.4},
  yticklabel style={/pgf/number format/fixed}]
\addplot[draw=none,name path=hiSEG0,forget plot] coordinates {(460,0.0782) (470,0.0782) (480,0.0811) (490,0.0844) (500,0.0861) (510,0.0894) (520,0.0934) (530,0.1010) (540,0.1081) (550,0.1119) (560,0.1135) (570,0.1150) (580,0.1173) (590,0.1168) (600,0.1186) (610,0.1202) (620,0.1214) (630,0.1232) (640,0.1245) (650,0.1246) (660,0.1251) (670,0.1267) (680,0.1290) (690,0.1358) (700,0.1468) (710,0.1622) (720,0.1824) (730,0.2057) (740,0.2273) (750,0.2444) (760,0.2506) (770,0.2618) (780,0.2684) (790,0.2673) (800,0.2732) (810,0.2780) (820,0.2762) (830,0.2818) (840,0.2861) (850,0.2881) (860,0.2903) (870,0.2921) (880,0.2991) (890,0.3027) (900,0.3011) (910,0.2995)};
\addplot[draw=none,name path=loSEG0,forget plot] coordinates {(460,0.0078) (470,0.0078) (480,0.0100) (490,0.0117) (500,0.0125) (510,0.0149) (520,0.0194) (530,0.0250) (540,0.0310) (550,0.0341) (560,0.0349) (570,0.0340) (580,0.0325) (590,0.0314) (600,0.0314) (610,0.0310) (620,0.0303) (630,0.0302) (640,0.0298) (650,0.0285) (660,0.0266) (670,0.0255) (680,0.0264) (690,0.0331) (700,0.0488) (710,0.0676) (720,0.0862) (730,0.1002) (740,0.1082) (750,0.1144) (760,0.1157) (770,0.1208) (780,0.1241) (790,0.1237) (800,0.1263) (810,0.1285) (820,0.1291) (830,0.1314) (840,0.1322) (850,0.1329) (860,0.1338) (870,0.1346) (880,0.1375) (890,0.1388) (900,0.1386) (910,0.1372)};
\addplot[spec!25,forget plot] fill between[of=hiSEG0 and loSEG0];
\addplot[spec,thick,mark=none,forget plot] coordinates {(460,0.0430) (470,0.0430) (480,0.0455) (490,0.0481) (500,0.0493) (510,0.0521) (520,0.0564) (530,0.0630) (540,0.0696) (550,0.0730) (560,0.0742) (570,0.0745) (580,0.0749) (590,0.0741) (600,0.0750) (610,0.0756) (620,0.0758) (630,0.0767) (640,0.0771) (650,0.0766) (660,0.0758) (670,0.0761) (680,0.0777) (690,0.0845) (700,0.0978) (710,0.1149) (720,0.1343) (730,0.1530) (740,0.1677) (750,0.1794) (760,0.1831) (770,0.1913) (780,0.1963) (790,0.1955) (800,0.1997) (810,0.2032) (820,0.2026) (830,0.2066) (840,0.2092) (850,0.2105) (860,0.2120) (870,0.2133) (880,0.2183) (890,0.2207) (900,0.2199) (910,0.2183)};
\addplot[draw=none,name path=hiSEG1,forget plot] coordinates {(980,0.3145) (990,0.3170) (1000,0.3211) (1010,0.3244) (1020,0.3297) (1030,0.3355) (1040,0.3392) (1050,0.3431) (1060,0.3477) (1070,0.3502) (1080,0.3502) (1090,0.3521) (1100,0.3507) (1110,0.3397) (1120,0.3127) (1130,0.3228)};
\addplot[draw=none,name path=loSEG1,forget plot] coordinates {(980,0.1513) (990,0.1516) (1000,0.1527) (1010,0.1539) (1020,0.1559) (1030,0.1580) (1040,0.1590) (1050,0.1603) (1060,0.1623) (1070,0.1637) (1080,0.1638) (1090,0.1650) (1100,0.1654) (1110,0.1610) (1120,0.1404) (1130,0.1446)};
\addplot[spec!25,forget plot] fill between[of=hiSEG1 and loSEG1];
\addplot[spec,thick,mark=none,forget plot] coordinates {(980,0.2329) (990,0.2343) (1000,0.2369) (1010,0.2392) (1020,0.2428) (1030,0.2468) (1040,0.2491) (1050,0.2517) (1060,0.2550) (1070,0.2569) (1080,0.2570) (1090,0.2586) (1100,0.2581) (1110,0.2504) (1120,0.2266) (1130,0.2337)};
\addplot[draw=none,name path=hiSEG2,forget plot] coordinates {(1160,0.3159) (1170,0.3196) (1180,0.3267) (1190,0.3358) (1200,0.3381) (1210,0.3375) (1220,0.3402) (1230,0.3433) (1240,0.3466) (1250,0.3565) (1260,0.3693) (1270,0.3779) (1280,0.3695) (1290,0.3563) (1300,0.3457) (1310,0.3395) (1320,0.3359) (1330,0.3270) (1340,0.3214)};
\addplot[draw=none,name path=loSEG2,forget plot] coordinates {(1160,0.1523) (1170,0.1544) (1180,0.1568) (1190,0.1598) (1200,0.1602) (1210,0.1602) (1220,0.1613) (1230,0.1625) (1240,0.1638) (1250,0.1675) (1260,0.1724) (1270,0.1771) (1280,0.1737) (1290,0.1687) (1300,0.1647) (1310,0.1625) (1320,0.1616) (1330,0.1572) (1340,0.1509)};
\addplot[spec!25,forget plot] fill between[of=hiSEG2 and loSEG2];
\addplot[spec,thick,mark=none,forget plot] coordinates {(1160,0.2341) (1170,0.2370) (1180,0.2418) (1190,0.2478) (1200,0.2491) (1210,0.2488) (1220,0.2508) (1230,0.2529) (1240,0.2552) (1250,0.2620) (1260,0.2708) (1270,0.2775) (1280,0.2716) (1290,0.2625) (1300,0.2552) (1310,0.2510) (1320,0.2487) (1330,0.2421) (1340,0.2361)};
\addplot[draw=none,name path=hiSEG3,forget plot] coordinates {(1460,0.2286) (1470,0.2073) (1480,0.2057) (1490,0.2212) (1500,0.2320) (1510,0.2432) (1520,0.2529) (1530,0.2572) (1540,0.2604) (1550,0.2645) (1560,0.2685) (1570,0.2729) (1580,0.2784) (1590,0.2829) (1600,0.2856) (1610,0.2888) (1620,0.2905) (1630,0.2913) (1640,0.2925) (1650,0.2939) (1660,0.2944) (1670,0.2940) (1680,0.2933) (1690,0.2901) (1700,0.2875) (1710,0.2855) (1720,0.2838) (1730,0.2808) (1740,0.2779) (1750,0.2751) (1760,0.2708) (1770,0.2627) (1780,0.2560) (1790,0.2506)};
\addplot[draw=none,name path=loSEG3,forget plot] coordinates {(1460,0.0676) (1470,0.0611) (1480,0.0625) (1490,0.0721) (1500,0.0776) (1510,0.0824) (1520,0.0869) (1530,0.0901) (1540,0.0932) (1550,0.0965) (1560,0.0996) (1570,0.1028) (1580,0.1061) (1590,0.1088) (1600,0.1111) (1610,0.1139) (1620,0.1156) (1630,0.1166) (1640,0.1177) (1650,0.1188) (1660,0.1190) (1670,0.1189) (1680,0.1185) (1690,0.1172) (1700,0.1158) (1710,0.1143) (1720,0.1127) (1730,0.1111) (1740,0.1100) (1750,0.1087) (1760,0.1063) (1770,0.1040) (1780,0.1012) (1790,0.0967)};
\addplot[spec!25,forget plot] fill between[of=hiSEG3 and loSEG3];
\addplot[spec,thick,mark=none,forget plot] coordinates {(1460,0.1481) (1470,0.1342) (1480,0.1341) (1490,0.1467) (1500,0.1548) (1510,0.1628) (1520,0.1699) (1530,0.1737) (1540,0.1768) (1550,0.1805) (1560,0.1840) (1570,0.1878) (1580,0.1923) (1590,0.1959) (1600,0.1984) (1610,0.2013) (1620,0.2030) (1630,0.2039) (1640,0.2051) (1650,0.2064) (1660,0.2067) (1670,0.2064) (1680,0.2059) (1690,0.2037) (1700,0.2016) (1710,0.1999) (1720,0.1982) (1730,0.1960) (1740,0.1939) (1750,0.1919) (1760,0.1886) (1770,0.1833) (1780,0.1786) (1790,0.1736)};
\addplot[draw=none,name path=hiSEG4,forget plot] coordinates {(1960,0.1574) (1970,0.1758) (1980,0.1902) (1990,0.2043) (2000,0.1809) (2010,0.1734) (2020,0.1755) (2030,0.2086) (2040,0.2199) (2050,0.2139) (2060,0.2215) (2070,0.2099) (2080,0.2152) (2090,0.2117) (2100,0.2099) (2110,0.2098) (2120,0.2088) (2130,0.2100) (2140,0.2125) (2150,0.2110) (2160,0.2098) (2170,0.2136) (2180,0.2122) (2190,0.2153) (2200,0.2114) (2210,0.2169) (2220,0.2166) (2230,0.2150) (2240,0.2128) (2250,0.2074) (2260,0.2041) (2270,0.2005) (2280,0.1980) (2290,0.1941) (2300,0.1886) (2310,0.1794) (2320,0.1667) (2330,0.1766) (2340,0.1729) (2350,0.1594) (2360,0.1408) (2370,0.1317) (2380,0.1278) (2390,0.1284) (2400,0.1304) (2410,0.1198) (2420,0.1160) (2430,0.1248) (2440,0.1251) (2450,0.1224)};
\addplot[draw=none,name path=loSEG4,forget plot] coordinates {(1960,0.0334) (1970,0.0368) (1980,0.0398) (1990,0.0414) (2000,0.0377) (2010,0.0386) (2020,0.0414) (2030,0.0480) (2040,0.0502) (2050,0.0491) (2060,0.0510) (2070,0.0520) (2080,0.0538) (2090,0.0540) (2100,0.0549) (2110,0.0559) (2120,0.0565) (2130,0.0574) (2140,0.0583) (2150,0.0590) (2160,0.0599) (2170,0.0611) (2180,0.0619) (2190,0.0634) (2200,0.0634) (2210,0.0650) (2220,0.0647) (2230,0.0638) (2240,0.0620) (2250,0.0594) (2260,0.0569) (2270,0.0545) (2280,0.0528) (2290,0.0507) (2300,0.0486) (2310,0.0464) (2320,0.0447) (2330,0.0453) (2340,0.0436) (2350,0.0413) (2360,0.0348) (2370,0.0301) (2380,0.0259) (2390,0.0257) (2400,0.0269) (2410,0.0199) (2420,0.0153) (2430,0.0195) (2440,0.0146) (2450,0.0016)};
\addplot[spec!25,forget plot] fill between[of=hiSEG4 and loSEG4];
\addplot[spec,thick,mark=none,forget plot] coordinates {(1960,0.0954) (1970,0.1063) (1980,0.1150) (1990,0.1228) (2000,0.1093) (2010,0.1060) (2020,0.1085) (2030,0.1283) (2040,0.1351) (2050,0.1315) (2060,0.1362) (2070,0.1309) (2080,0.1345) (2090,0.1329) (2100,0.1324) (2110,0.1329) (2120,0.1327) (2130,0.1337) (2140,0.1354) (2150,0.1350) (2160,0.1349) (2170,0.1374) (2180,0.1371) (2190,0.1394) (2200,0.1374) (2210,0.1410) (2220,0.1407) (2230,0.1394) (2240,0.1374) (2250,0.1334) (2260,0.1305) (2270,0.1275) (2280,0.1254) (2290,0.1224) (2300,0.1186) (2310,0.1129) (2320,0.1057) (2330,0.1109) (2340,0.1083) (2350,0.1004) (2360,0.0878) (2370,0.0809) (2380,0.0769) (2390,0.0770) (2400,0.0787) (2410,0.0698) (2420,0.0656) (2430,0.0722) (2440,0.0699) (2450,0.0620)};
\node[font=\footnotesize,anchor=south west] at (axis cs:560,0.29) {red edge};
\draw[->,attn] (axis cs:640,0.29) -- (axis cs:715,0.14);
\end{axis}
\end{tikzpicture}
\caption{The mean reflectance of each of the $165$ common-grid bands over the normalisation corpus
(teal line), with one standard deviation either side (shaded). Real numbers from
\texttt{app/constants/hyperspectral.json}, plotted at the grid wavelengths; the four breaks are
the atmospheric windows removed from the grid.}\label{fig:light:meanspec}
\end{figure}

\section{The cube}\label{sec:light:cube}

Stack one image per band and you get a three-dimensional block of numbers.

\begin{definition}[Array, tensor and shape]\index{array}\index{tensor}\index{shape}
An \textbf{array} is a block of numbers indexed by one or more whole-number positions. Its
\textbf{shape} lists how many positions each index has. A one-index array is a vector, a
two-index array is a matrix, and arrays with more indices are commonly called
\textbf{tensors}. In this book ``tensor'' means a multi-index array of numbers, nothing more.
\end{definition}

\begin{definition}[Hyperspectral cube]\index{hyperspectral cube}
A \textbf{hyperspectral cube} is an array $\mat{X}$ of shape \shape{(C, H, W)}: $C$ bands,
$H$ rows, $W$ columns. Its entry $\mat{X}[c, i, j] = x_c(i,j)$ is the reflectance of band $c$ at
pixel $(i,j)$. A sensor that records $200$ or more narrow bands is called
\textbf{hyperspectral}\index{hyperspectral}.
\end{definition}

``Hyperspectral'' means high \emph{spectral} resolution: many narrow wavelength ranges. It does
not mean sharp images. Spatially these sensors are coarse; a PRISMA or EnMAP pixel covers about
$30$ metres on the ground. The ``cube'' is not three-dimensional on the ground either: every
layer is a picture of the same flat ground, and the third axis is wavelength, not height.

The cube can be read in two ways, both correct and both used.
\begin{itemize}
\item \textbf{A stack of images.} Fix the band $c$ and take $\mat{X}[c, :, :]$, an
\shape{(H, W)} array: the whole scene at one wavelength, a grey picture. Visualisation uses this
reading.
\item \textbf{A grid of spectra.} Fix the pixel $(i,j)$ and take $\mat{X}[:, i, j]$, a vector of
length $C$: one pixel at every wavelength. Detectors use this reading.
\end{itemize}
In Python the colon ``\code{:}'' means ``every position along this axis''.

\coursefig[0.95\linewidth]{iw-02-cube}{From photograph to spectrum. (a) A colour photograph
stores three bands per pixel. (b) A hyperspectral cube stores one image per wavelength, with shape
\texttt{(C, H, W)}. (c) Reading one pixel down the band axis gives its spectrum. (Course figure,
Indradhanu part 2; the curve is illustrative.)}{fig:light:cube}

\begin{worked}[label={we:light:cube}]{Slicing a 3-band, 2-by-2 cube}
This small cube is used again in \cref{ch:bandpipeline}, so copy it down. It has $C=3$ bands, at
$460$, $860$ and $1{,}650$ nm, and $H = W = 2$ (values illustrative\tagA):
\[
\mat{X}[0] = \begin{pmatrix} 0.03 & 0.05 \\ 0.05 & 0.03 \end{pmatrix},\quad
\mat{X}[1] = \begin{pmatrix} 0.20 & 0.30 \\ 0.30 & 0.20 \end{pmatrix},\quad
\mat{X}[2] = \begin{pmatrix} 0.14 & 0.10 \\ 0.14 & 0.10 \end{pmatrix}.
\]
Each matrix is one band; rows are $i = 0, 1$ and columns $j = 0, 1$.

\emph{A band.} $\mat{X}[1, :, :]$ is the $860$ nm image, the middle matrix.

\emph{A spectrum.} For pixel $(0,1)$ (row $0$, column $1$) read position $(0,1)$ of each
matrix in turn:
\begin{align*}
\mat{X}[0,0,1] &= 0.05,\\
\mat{X}[1,0,1] &= 0.30,\\
\mat{X}[2,0,1] &= 0.10,
\end{align*}
so $\vect{x}_{01} = (0.05,\ 0.30,\ 0.10)$. The other three spectra are
$\vect{x}_{00} = (0.03,\ 0.20,\ 0.14)$, $\vect{x}_{10} = (0.05,\ 0.30,\ 0.14)$ and
$\vect{x}_{11} = (0.03,\ 0.20,\ 0.10)$.

\emph{Counting.} The cube holds $3 \times 2 \times 2 = 12$ numbers: $3$ bands of $4$ pixels, or
equivalently $4$ spectra of $3$ values.
\end{worked}

\begin{worked}[label={we:light:bytes}]{The size of a real cube}
A PRISMA cube has shape \shape{(239, 1210, 1219)} stored as \code{float32} (4 bytes per
number). How large is it? Count the pixels, then the numbers, then the bytes:
\begin{align*}
1{,}210 \times 1{,}219 &= 1{,}474{,}990 \text{ pixels},\\
1{,}474{,}990 \times 239 &= 352{,}522{,}610 \text{ numbers},\\
352{,}522{,}610 \times 4 &= 1{,}410{,}090{,}440 \text{ bytes} \approx 1.41\ \text{GB}.
\end{align*}
After resampling to the $165$-band grid of \cref{ch:sensors} the same scene needs
\[ 1{,}474{,}990 \times 165 \times 4 = 973{,}493{,}400 \text{ bytes} \approx 0.97\ \text{GB}. \]
\end{worked}

\begin{table}[htbp]\centering
\caption{Vocabulary to pin down now. The two easiest to confuse are band and spectrum: a band is
a horizontal slice (one wavelength, all pixels), a spectrum a vertical skewer (one pixel, all
wavelengths).}\label{tab:light:vocab}
\begin{tabular}{@{}lll@{}}\toprule
Word & What it is & Shape \\\midrule
band & one wavelength layer, a grey picture & \shape{(H, W)}\\
spectrum & one pixel's values across all bands & \shape{(C,)}\\
cube & the whole stack & \shape{(C, H, W)}\\
patch & a small square cut from a cube for training & \shape{(165, 128, 128)}\\
scene & one complete acquisition & \shape{(C, about 1000, about 1000)}\\
validity cube & the 0/1 companion of the cube (\cref{sec:light:validity}) & \shape{(C, H, W)}\\\bottomrule
\end{tabular}
\end{table}

\section{Thermal images: the one-band cousin}\label{sec:light:thermal}

Two of the five sensors, Landsat~9 and HotSat-1, are \textbf{thermal}\index{thermal sensor}: they
record one band in the thermal infrared, which measures how much heat the ground
\emph{emits}, and so how hot it is. A thermal image is the same cube structure with $C = 1$,
shape \shape{(1, H, W)}. ``Hotter than its surroundings'' is exactly what a fire, a gas flare or a
running machine looks like, so anomaly detection is still useful, but there is no spectrum to
compare. The thermal model Chakshu shares almost all its code with Indradhanu; everywhere this
book says ``$165$ bands'', Chakshu says ``$1$ band''.

A sensor does not store physics directly. It stores \textbf{digital numbers}\index{digital
number} (DN): integers proportional to the energy that hit the detector, scaled by the
electronics. DN are not comparable across sensors, dates or sun angles, so the first processing
step converts them to a physical quantity. For Landsat's surface-temperature product the
conversion is a straight line with two published constants\tagM{l2sp_dn_to_temperature_transformer.py},
in \file{app/utils/data_transformations/l2sp_dn_to_temperature_transformer.py}:
\begin{equation}\label{eq:light:landsat}
T_{\mathrm{K}} = 0.00341802 \times \mathrm{DN} + 149.0, \qquad T_{^\circ\mathrm{C}} = T_{\mathrm{K}} - 273.15.
\end{equation}
The first gives temperature in kelvin (K), the scale that starts at absolute zero; subtracting
$273.15$ converts to degrees Celsius.

\begin{worked}[label={we:light:landsat}]{A Landsat digital number to Celsius}
A Landsat pixel reports $\mathrm{DN} = 30{,}000$. Do each step on paper.

Multiply by the scale factor:
\[ 0.00341802 \times 30{,}000 = 102.5406. \]
Add the offset to get kelvin:
\[ T_{\mathrm{K}} = 102.5406 + 149.0 = 251.5406\ \mathrm{K}. \]
Subtract $273.15$:
\[ T_{^\circ\mathrm{C}} = 251.5406 - 273.15 = -21.6094 \approx -21.61\ ^\circ\mathrm{C}. \]
Minus $21.6$ degrees: plausibly the top of a cold cloud.
\end{worked}

\section{Validity: not every number is a measurement}\label{sec:light:validity}

\Cref{ch:problem} introduced the validity mask. For hyperspectral data the codebase keeps one
validity value per band per pixel.

\begin{definition}[Validity cube]\index{validity cube}
The \textbf{validity cube} $\mat{V}$ has the same shape \shape{(C, H, W)} as the cube, with
$\mat{V}[c,i,j] = 1$ if band $c$ of pixel $(i,j)$ is a real measurement and $0$ if it is not
(dead detector, cloud, off-swath).
\end{definition}

After the processing of \cref{ch:bandpipeline}, hyperspectral validity is \emph{all-or-nothing per
pixel}: a pixel is valid in all bands or in none. That lets code collapse the cube to one
\shape{(H, W)} map, and the codebase does it in two different ways. The training code reads band
$0$ alone (\file{app/foundation_models/trainers/hyperspectral_segformer_mae_trainer.py}, line 160):
\begin{lstlisting}
# Use band 0 as spatial validity proxy: (B, C, H, W) -> (B, 1, H, W)
return validity[:, 0:1, :, :].float()
\end{lstlisting}
The scoring code sums over bands and asks whether any band is valid
(\file{backend/allotrope/action_types/_anomaly_scoring_run.py}, line 174):
\begin{lstlisting}
spatial_valid = (validity_np.sum(axis=0) > 0).astype(np.uint8)   # (H, W)
\end{lstlisting}
Here \code{B} is the \textbf{batch}\index{batch} size, the number of image patches processed
together in one step; \code{sum(axis=0)} adds along the first axis (the bands).

\begin{worked}[label={we:light:validity}]{Two ways to collapse validity}
A $3$-band, $1 \times 3$ image (one row, three pixels) has validity
\[
\mat{V}[0] = (1,\ 0,\ 0),\quad \mat{V}[1] = (1,\ 0,\ 1),\quad \mat{V}[2] = (1,\ 0,\ 1).
\]
Sum over bands, pixel by pixel:
\begin{align*}
\text{pixel } 0:\ 1 + 1 + 1 &= 3,\\
\text{pixel } 1:\ 0 + 0 + 0 &= 0,\\
\text{pixel } 2:\ 0 + 1 + 1 &= 2.
\end{align*}
``Sum $> 0$'' gives $(1,\ 0,\ 1)$. ``Band $0$ only'' gives $(1,\ 0,\ 0)$. They disagree on pixel
$2$, which is valid in two bands but not in band $0$. That pixel is \emph{not} all-or-nothing, so
this validity cube has not been through the pipeline. After the pipeline, pixel $2$ would have been
either filled to $(1,1,1)$ or wiped to $(0,0,0)$ (\cref{ch:bandpipeline}), and both methods would
agree. The training code's shortcut reads one band instead of $165$, which is $165$ times less
memory traffic on an operation done for every batch; it is safe only because of the
all-or-nothing guarantee.
\end{worked}

\begin{keypoint}
The two collapses agree only on cubes that have been through the band pipeline; on raw validity
the band-0 shortcut is wrong.
\end{keypoint}

\begin{confusions}
\confusion{Is a band the same as a channel?}{In practice yes. Remote sensing says band, deep
learning says channel; the $C$ in \shape{(C, H, W)} is the band count either way.}
\confusion{Does hyperspectral mean high resolution?}{It means high \emph{spectral} resolution,
many narrow bands. Spatially a pixel is about $30$ m, coarse by photographic standards.}
\confusion{Is the band value the reflectance at the centre wavelength?}{Not exactly. It is a
weighted average over the band's response function (\cref{eq:light:srf}). For narrow bands and a
smooth spectrum the two are close; near a sharp feature they differ
(\cref{we:light:srf}).}
\confusion{Are more bands always better?}{No. Adjacent bands $10$ nm apart carry nearly the same
information, each narrow band is noisier, and many bands relative to usable pixels make some
statistics numerically unstable (\cref{ch:rx}). Indradhanu's first layer compresses the $165$
bands for exactly this reason (\cref{ch:compressor}).}
\confusion{Reflectance is between $0$ and $1$, so values above $1$ are bugs?}{Usually, not always.
Specular glint off water or metal can exceed $1$. Do not clamp blindly.}
\confusion{Radiance or reflectance: which does the model see?}{Reflectance. The providers' level-2
products are already surface reflectance (as scaled integers); the codebase undoes the integer
scaling and never handles radiance.}
\confusion{If validity is all-or-nothing, why is the validity cube $C$ layers deep?}{Because the
same type is used at earlier pipeline stages where validity genuinely varies per band. By the time
a model sees it, the layers agree.}
\end{confusions}

\begin{checkyourself}
\item\label{ex:light:1} A band has FWHM $12$ nm. Compute $\sigma$ to three decimal places, then the
Gaussian weight $3$ nm from the centre. Show every step.
\item\label{ex:light:2} A band samples wavelengths $1{,}000$, $1{,}005$, $1{,}010$ nm with weights
$0.5$, $1$, $0.5$; the reflectances there are $0.30$, $0.26$, $0.18$. Compute the band value and
compare it with the reflectance at the centre.
\item\label{ex:light:3} A Landsat pixel reports $\mathrm{DN} = 41{,}000$. Compute its temperature in
kelvin and in Celsius.
\item\label{ex:light:4} Using the cube of \cref{we:light:cube}, write the spectrum of pixel
$(1,0)$ and the $1{,}650$ nm band as a matrix. Is pixel $(1,0)$ a scaled copy of pixel $(0,0)$?
Check by dividing component by component.
\item\label{ex:light:5} With the numbers of \cref{we:light:rho}, a second surface gives radiance
$L = 0.05$ under the same sun. What is its reflectance? Then explain in one or two sentences why
reflectance, not radiance, is what lets one model be used on several sensors.
\end{checkyourself}

\nextchapter{ch:sensors}
